\documentclass[12pt]{extarticle}
\usepackage{exam}
\begin{document}
\begin{coverpage}
	{\fontsize{13pt}{13pt}\selectfont W\textsuperscript{2} Education Standards Authority}
	\vspace{2.4cm}

	\examnameyear{2077}{Highest School Certificate Examination}

	\coursename{33pt}{Mathematics Extension 1}

	\vspace{1cm}

		\begin{generalinstructions}
			\item Reading time -- 10 minutes
			\item Working time -- 2 hours
			\item Write using black pen
			\item Calculators approved by NESA may be used
			\item A reference sheet is provided at the back of this paper
			\item For questions in Section II, show relevant mathematical reasoning and/or calculations
			\item Write your Centre Number and Student Number at the top of \\ this page
		\end{generalinstructions}

	\vspace{1cm}

	\begin{sectionlist}{70}
		\textbf{Section I -- 10 marks} (pages 2--6)
		\begin{itemize}
			\item Attempt Questions 1--10
			\item Allow about 15 minutes for this section
		\end{itemize}
		\textbf{Section II -- 60 marks} (pages 7--12)
		\begin{itemize}
			\item Attempt Questions 11--14
			\item Allow about 1 hour and 45 minutes for this section
		\end{itemize}
	\end{sectionlist}
\end{coverpage}


\begin{sectionheader}{Section I}
	\textbf{10 marks\\
		Attempt Questions 1--10\\
		Allow about 15 minutes for this section}

	Use the multiple-choice answer sheet for Questions 1--10.
\end{sectionheader}
\begin{multiplechoice}
	\item Given $\tan\theta=\frac{1}{2}$, what is the exact value of $\tan\qty(\theta-\frac{\pi}{6})$?
	\begin{enumerate}
		\item $8+5\sqrt{3}$
		\item $8-5\sqrt{3}$
		\item $\dfrac{8+5\sqrt{3}}{11}$
		\item $\dfrac{8-5\sqrt{3}}{11}$
	\end{enumerate}

	\item The function $f$ is defined by the follwing parametric equations:
	\[x = 2e^t,\, y = \cos(1+e^{3t}).\]
	Which of the following is equivalent to $\dv{y}{x}$?
	\begin{enumerate}
		\item $-\frac{3}{2}e^{2t}\sin(1+e^{3t})$
		\item $\frac{3}{2}e^{2t}\sin(1+e^{3t})$
		\item $6e^{2t}\sin(1+e^{3t})$
		\item $-6e^{2t}\sin(1+e^{3t})$
	\end{enumerate}
	\pagebreak
	\item Daniel has recently been into collecting Labubus and wants to display a total of five different Labubus on his shelf.
	He can choose between five different Labubus from the \textit{Have a Seat} collection and five different Labubus from the \textit{Exiting Macaron} collection.

	Given that the pink \textit{Exiting Macaron} Labubu must be displayed, how many ways can he select the remaining if he wants at least two from the \textit{Have a Seat} collection and at least one other from the \textit{Exiting Macaron} collection on the shelf?
	\begin{enumerate}
		\item $100$
		\item $240$
		\item $252$
		\item $320$
	\end{enumerate}
	\item Which of the following differential equations matches the slope field shown?
	% \begin{noindent}
  \begin{center}
  \begin{asy}
    import graph;
    size(7cm);
    real f(real x, real y) {
      return 1/(y-x);
    }
    real xmin = -10, xmax = 10;
    real ymin = -10, ymax = 10;
    int nx = 20, ny = 20;
    real len = 0.35;
    for (int i = 0; i <= nx; ++i) {
      real x = xmin + (xmax - xmin) * i / nx;
      for (int j = 0; j <= ny; ++j) {
        real y = ymin + (ymax - ymin) * j / ny;
        if (x == y) {
          continue;
        }
        real m = f(x, y);
        real dx = len / sqrt(1 + m^2);
        real dy = m * dx;
        draw((x - dx, y - dy) -- (x + dx, y + dy),
             linewidth(0.7));
      }
    }
    draw((xmin, 0) -- (xmax, 0), Arrow);
    draw((0, ymin) -- (0, ymax), Arrow);
    label("$x$", (xmax+0.5, 0), E);
    label("$y$", (0, ymax+0.5), N);
  \end{asy}
  \end{center}
  % \end{noindent}
	\begin{enumerate}
		\item $\dv{y}{x}=y-x$
		\item $\dv{y}{x}=x-y$
		\item $\dv{y}{x}=\frac{1}{y-x}$
		\item $\dv{y}{x}=\frac{1}{x-y}$
	\end{enumerate}
	\pagebreak
	\item The polynomial $P(x)=6x^4-x^3-17x^2+16x-4$ has four real roots.
	What is the sum of the roots of the polynomial $P\qty(-\frac{1}{2}x-1)$?
	\begin{enumerate}
		\item $\frac{5}{6}$
		\item $\frac{1}{12}$
		\item $-2\frac{1}{3}$
		\item $-8\frac{1}{3}$
	\end{enumerate}
	\item What is the constant term in the binomial expansion of $\qty(4x^2-\frac{1}{x})^9$?
	\begin{enumerate}
		\item $-344064$
		\item $-5376$
		\item $84$
		\item $5376$
	\end{enumerate}
	\item Which of the following expressions is eqivalent to $\int\frac{4}{x}\qty(\ln x)^3\dd{x}$?
	\begin{enumerate}
		\item $\qty(\ln x)^4 + C$
		\item $4x^4\ln x + C$
		\item $\frac{4x^4}{\ln x} + C$
		\item $\frac{4}{x^4} + C$
	\end{enumerate}
	\pagebreak
	\item In the diagram below, $AOB$ is a diameter of the circle with centre $O$.
	Point $C$ lies on the circumference of the circle.
	% \begin{noindent}
  \begin{center}
  \begin{asy}
    size(7cm);
    draw(circle((0,0), 1), black);
    draw((-1,0) -- (1,0), black);
    draw((0,0) -- (1/2, sqrt(3)/2), black);
    draw((1,0) -- (1/2, sqrt(3)/2), black);
    draw((-1,0) -- (1/2, sqrt(3)/2), black);
    label("$A$", (-1,0), W);
    label("$B$", (1,0), E);
    label("$O$", (0,0), S);
    label("$C$", (1/2, sqrt(3)/2), NE);
  \end{asy}
  \end{center}
  % \end{noindent}
	If $\overrightarrow{OC}=\mathbf{r}$ and $\overrightarrow{BC}=\mathbf{s}$, to which of the following is $\overrightarrow{AC}$ equal?
	\begin{enumerate}
		\item $\mathbf{r}+2\mathbf{s}$
		\item $\mathbf{r}-2\mathbf{s}$
		\item $2\mathbf{r}+\mathbf{s}$
		\item $2\mathbf{r}-\mathbf{s}$
	\end{enumerate}
	\item A particle is projected from a window 9 metres above the ground, at an angle of $\theta$ above the horizontal, where $\tan\theta=\frac{3}{4}$ and the initial velocity is $20ms^{-1}$.

	Taking $g=10$, what is the maximum height reached by the particle?
	\begin{enumerate}
		\item $6.5$ metres
		\item $7.2$ metres
		\item $15.5$ metres
		\item $16.2$ metres
	\end{enumerate}
	\pagebreak
	\item The graph of $y=2\cos(\frac{\pi x}{3})+1$ is shown.
	% \begin{noindent}
  \begin{center}
  \begin{asy}
    import graph;
    size(10cm);
    real f(real x) {
      return 2*cos(pi*x/3)+1;
    }
    path g = graph(f, 0, 3);
    path ga = graph(f, 0, 2);
    path gb = graph(f, 2, 3);
    draw(g, black);
    draw((-1.5,0)--(3.5,0), black, arrow=Arrows);
    draw((0,-1.5)--(0,3.5), black, arrow=Arrows);
    label("$O$", (0,0), SW);
    label("$y$", (0,3.5), N);
    label("$x$", (3.5,0), E);
    fill((0,3)--ga--(0,0)--cycle, grey + opacity(0.3));
    fill((2,0)--gb--(3,0)--cycle, grey + opacity(0.3));
    label("$A$", (0.7,1.25));
    label("$B$", (2.65,-0.45));
    for (int i = -1; i <= 3; ++i) {
      draw((i,-0.08)--(i,0.08), black);
      if (i != 0)
        label(string(i), (i,0), 3S);
    }
    for (int i = -1; i <= 3; ++i) {
      draw((-0.08,i)--(0.08,i), black);
      if (i != 0)
        label(string(i), (0,i), 3W);
    }
  \end{asy}
  \end{center}
  % \end{noindent}
	Region $A$ has an area of $\frac{3\sqrt{3}}{\pi}+2$ square units.

	Region $B$ has an area of $\frac{3\sqrt{3}}{\pi}-1$ square units.

	Using this information, what is the value of $\int^{3}_{-1}\frac{3}{\pi}\cos^{-1}\qty(\frac{x-1}{2})\dd{x}$?
	\begin{enumerate}
		\item $3$
		\item $6$
		\item $\frac{6\sqrt{3}}{\pi}+1$
		\item $11-\frac{6\sqrt{3}}{\pi}$
	\end{enumerate}
\end{multiplechoice}
\pagebreak
\begin{sectionheader}{Section II}
	\textbf{60 marks\\
		Attempt Questions 11--14\\
		Allow about 1 hour and 45 minutes for this section}

	Answer each question in the appropriate writing booklet. Extra writing booklets are available.

	For questions in Section II, your responses should include relevant mathematical reasoning and/or calculations.
\end{sectionheader}
\begin{shortanswer}
	\begin{bookletquestion}{12}
		\begin{enumerate}
			\item \mrks{2} Consider the vectors $\mathbf{a}=m\boldsymbol{\hat{\imath}}+\boldsymbol{\hat{\jmath}}$ and $\mathbf{b}=\boldsymbol{\hat{\imath}}+m\boldsymbol{\hat{\jmath}}$, where $m\in\mathbb{R}$.
			      Find the values of $m$ for which the acute angle between $\mathbf{a}$ and $\mathbf{b}$ is $\frac{\pi}{3}$.
			\item The letters from the word VOLUME are placed at random on the circumference of a circle without replacement.
			      \begin{enumerate}
				      \item \mrks{1} Find the total number of distinct arrangements.
			      \end{enumerate}
			      If one of the arrangements are chosen at random, find the probability that:
			      \begin{enumerate}[resume]
				      \item \mrks{2} all the vowels are together.
				      \item \mrks{1} the vowels and consonants alternate.
			      \end{enumerate}
			\item Given $P(x)=3x^3-2x^2+x-3$ has zeros $\alpha$, $\beta$ and $\gamma$:
			      \begin{enumerate}
				      \item \mrks{1} write the value of $\alpha\beta\gamma$.
				      \item \mrks{1} Hence, or otherwise, find the value of $\frac{1}{\alpha}+\frac{1}{\beta}+\frac{1}{\gamma}$.
			      \end{enumerate}
			\item \mrks{4} Prove by mathematical induction that, for all positive integers $n$,
			      \[\dfrac{2}{1\times 3}+\dfrac{2}{2\times 4}+\dfrac{2}{3\times 5}+\cdots+\dfrac{2}{n(n+2)}=\dfrac{3}{2}-\dfrac{2n+3}{(n+1)(n+2)}.\]
		\end{enumerate}
		\questionend
	\end{bookletquestion}
	\pagebreak
	\begin{bookletquestion}{??}
		\begin{enumerate}
			\item \begin{enumerate}
				      \item \mrks{2} Express $\sin x -\sqrt{3}\cos x$ in the form $R\sin(x-\alpha)$, where $R>0$ and \\ $0\leq\alpha\leq\frac{\pi}{2}$.
				      \item \mrks{1} Hence, solve $\sin x -\sqrt{3}\cos x=2$ for $0\leq x \leq \pi$.
			      \end{enumerate}
			\item Consider the differential equation $\dv{y}{x}=-(x^2-1)(y-1)$ and its slope field, which is graphed below.
			      % \begin{noindent}
            \begin{center}
            \begin{asy}
              import graph;
              size(7cm);
              real f(real x, real y) {
                return -(x^2-1)*(y-1);
              }
              real xmin = -2, xmax = 2;
              real ymin = -2, ymax = 2;
              int nx = 8, ny = 8;
              real len = 0.08;
              for (int i = 0; i <= nx; ++i) {
                real x = xmin + (xmax - xmin) * i / nx;
                for (int j = 0; j <= ny; ++j) {
                  real y = ymin + (ymax - ymin) * j / ny;
                  real m = f(x, y);
                  real dx = len / sqrt(1 + m^2);
                  real dy = m * dx;
                  draw((x - dx, y - dy) -- (x + dx, y + dy),
                       linewidth(0.7));
                }
              }
              draw((xmin-0.5, 0) -- (xmax+0.5, 0), Arrow);
              draw((0, ymin-0.5) -- (0, ymax+0.5), Arrow);
              label("$x$", (xmax+0.7, 0), E);
              label("$y$", (0, ymax+0.7), N);
              for (int i = -2; i <= 2; ++i) {
                if (i != 0)
                  label(string(i), (i,0), 2S);
              }
              for (int i = -2; i <= 2; ++i) {
                if (i != 0)
                  label(string(i), (0,i), 2W);
              }
            \end{asy}
            \end{center}
            % \end{noindent}
			      \begin{enumerate}
				      \item \mrks{1} By referring to the slope field for the differential equation, explain why a solution could not have the graph shown below.
                 
				            \hspace*{-\leftmargin}%
				            \begin{minipage}{\dimexpr\linewidth+\leftmargin\relax}
					            % \begin{noindent}
                      \begin{center}
                      \begin{asy}
                        import graph;
                        size(7cm);
                        real f(real x) {
                          return 1/2*(x-7/4)*(x+7/4)*x;
                        }
                        real xmin = -2.5, xmax = 2.5;
                        real ymin = -2.5, ymax = 2.5;
                        draw((xmin, 0) -- (xmax, 0), Arrow);
                        draw((0, ymin) -- (0, ymax), Arrow);
                        draw(graph(f, -2.2, 2.2), black, Arrows);
                        label("$x$", (xmax+0.2, 0), E);
                        label("$y$", (0, ymax+0.2), N);
                        for (int i = -2; i <= 2; ++i) {
                          draw((i,-0.08)--(i,0.08), black);
                          if (i != 0)
                            label(string(i), (i,0), 2S);
                        }
                        for (int i = -2; i <= 2; ++i) {
                          draw((-0.08,i)--(0.08,i), black);
                          if (i != 0)
                            label(string(i), (0,i), 2W);
                        }
                      \end{asy}
                      \end{center}
                      % \end{noindent}
				            \end{minipage}
				            \questionbreak
				      \item \mrks{3} Find the particular solution of the differential equation that passes through the point $(0, 2)$.
			      \end{enumerate}
		\end{enumerate}
		\questionend
		\centerbold{End of paper}
	\end{bookletquestion}
\end{shortanswer}
\pagebreak
\blankpage
\end{document}
